\NSPage{139}%
\begin{EESegment}{EE-TGT-P139-001}
At large \NSi{NS-F-P139-001}{X}, the heat profile approaches the original power law, with the derivative bounds in
\eqref{eq:A.38}. Let \NSi{NS-F-P139-002}{E_{\mathrm{cl}}} stand for the reference profile built in
Proposition~\ref{prop:A.4}. The terminal profile will be replaced by this small change, and earlier corrections will
restore the three moments it changes. The other two moments do not change, because \NSi{NS-F-P139-003}{U=0} on both
regions that are edited.
\end{EESegment}

\begin{proposition}\label{prop:A.7}
\begin{EESegment}{EE-TGT-P139-002}
Let \NSi{NS-F-P139-004}{X_K=e^{.2}X_{\mathrm{tail}}}. Choose a smooth step
\NSi{NS-F-P139-005}{0\leq\chi_K(y)\leq 1} that is zero for
\NSi{NS-F-P139-006}{y\leq .2} and one for \NSi{NS-F-P139-007}{y\geq .5}. Make the terminal replacement
\end{EESegment}
\NSFormula{NS-F-P139-008}{display}{A.39}
\begin{equation}
  E_{\mathrm{cl}}\longmapsto E_{\mathrm{cl}}
  \bigl[1+\chi_K(y)(\mathcal H(2d/X)-1)\bigr]
  \tag{A.39}\label{eq:A.39}
\end{equation}
\begin{EESegment}{EE-TGT-P139-003}
Three additive \NSi{NS-F-P139-009}{E}-bumps in the second reserved patch of the intermediate power-law interval can
compensate for this replacement. The full edit preserves the functions \NSi{NS-F-P139-010}{M,J}, the totals
\NSi{NS-F-P139-011}{S(\infty),C_p(\infty)} from \eqref{eq:4.15}, and the renormalized angular moment
\NSi{NS-F-P139-012}{\int_0^\infty(H-H_{\mathrm{pow}})\,dX}, at every point of
\NSi{NS-F-P139-013}{[-1,1]}. When \NSi{NS-F-P139-014}{X_R} is large enough, the edit also preserves
\NSi{NS-F-P139-015}{E>0} and the strict admissible stress cone from that patch through
\NSi{NS-F-P139-016}{y=.5}. It leaves both the original pressure datum on the axis and the profiles before the
compensation patch unchanged.
\end{EESegment}
\end{proposition}

\begin{proof}
\begin{EESegment}{EE-TGT-P139-004}
First estimate the three changes made by the heat factor. Then cancel them with three earlier angular bumps. The
normalized radial moment formulas will show that these corrections keep the stress inside the admissible cone.
\end{EESegment}

\begin{EESegment}{EE-TGT-P139-005}
Put \NSi{NS-F-P139-017}{x=X/X_K}, and let \NSi{NS-F-P139-018}{e_K} be the reference amplitude at
\NSi{NS-F-P139-019}{X_K}. The reference tail is comparable to \NSi{NS-F-P139-020}{e_Kx^{-A}}.
Equations~\eqref{eq:A.34}--\eqref{eq:A.36}, together with the chain rule, show that for every pair of fixed integers
\NSi{NS-F-P139-021}{j,m\geq 0},
\end{EESegment}
\NSFormula{NS-F-P139-022}{display}{none}
\[
  \bigl|(X\partial_X)^j\partial_\eta^m(E-E_{\mathrm{cl}})\bigr|
  \leq C_{j,m}e_KX_K^{-1}x^{-A-1},\qquad x\geq 1.
\]
\begin{EESegment}{EE-TGT-P139-006}
The constants may depend on the schedule, which is already fixed. In particular, every derivative in this estimate is
bounded at \NSi{NS-F-P139-023}{d=0}. Let \NSi{NS-F-P139-024}{\Delta} mean the change caused only by
\eqref{eq:A.39}. The three differences of moments then obey
\end{EESegment}
\NSFormula{NS-F-P139-025}{display}{A.40}
\begin{equation}
  \left|\partial_\eta^m\Delta\int_0^\infty\frac{E^2}{2X}\,dX\right|
  \leq C_m\frac{e_K^2}{X_K}\int_1^\infty x^{-2A-2}\,dx,
  \tag{A.40}\label{eq:A.40}
\end{equation}
\NSFormula{NS-F-P139-026}{display}{A.41}
\begin{equation}
  \left|\partial_\eta^m\Delta S(\infty)\right|
  \leq C_me_K^2\int_1^\infty x^{-2A-1}\,dx,
  \tag{A.41}\label{eq:A.41}
\end{equation}
\NSFormula{NS-F-P139-027}{display}{A.42}
\begin{equation}
  \left|\partial_\eta^m\Delta\int_0^\infty(H-H_{\mathrm{pow}})\,dX\right|
  \leq C_me_KX_K^{1/2}\int_1^\infty x^{-A-1/2}\,dx
  \leq C_{m,h}e_KX_K^{1/2}.
  \tag{A.42}\label{eq:A.42}
\end{equation}
\begin{EESegment}{EE-TGT-P139-007}
Here \NSi{NS-F-P139-028}{H=\sqrt{2X}E} and
\NSi{NS-F-P139-029}{H_{\mathrm{pow}}=\sqrt{2X}E_{\mathrm{pow}}}; the pure power itself is not changed. The last
integral equals \NSi{NS-F-P139-030}{1/h}. It is finite because
\NSi{NS-F-P139-031}{h>0} was fixed before the large radial scale was chosen.
\end{EESegment}

\begin{EESegment}{EE-TGT-P139-008}
Once each one is divided by its natural radial scale, the three differences of moments are small. Cancel them on the
second reserved patch, where the power-law profile supplies three independent weights for the moments. Let
\NSi{NS-F-P139-032}{X_*} be the start of this patch, and put \NSi{NS-F-P139-033}{x_*=X/X_*}. Before the correction,
the profile there is \NSi{NS-F-P139-034}{E=e_*f(\eta)x_*^{-1/2-\lambda}}, with
\NSi{NS-F-P139-035}{U=0} and \NSi{NS-F-P139-036}{f=(1+\eta^2)^{-1}\geq 1/2}. Choose three nonnegative bumps
\NSi{NS-F-P139-037}{\beta_i} whose supports are disjoint, ordered, and contained in this patch. Add
\end{EESegment}
\NSFormula{NS-F-P139-038}{display}{none}
\[
  \delta E=e_*\sum_{i=1}^3c_i(\eta)\beta_i(x_*).
\]
\begin{EESegment}{EE-TGT-P139-009}
Normalize the pressure row, the \NSi{NS-F-P139-039}{S} row, and the angular row, in that order, by
\end{EESegment}
\NSFormula{NS-F-P139-040}{display}{A.43}
\begin{equation}
  e_*^2,\qquad X_*e_*^2,\qquad X_*^{3/2}e_*.
  \tag{A.43}\label{eq:A.43}
\end{equation}
\begin{EESegment}{EE-TGT-P139-010}
At \NSi{NS-F-P139-041}{c=0}, their derivatives are the integrals against the weights
\end{EESegment}
\NSFormula{NS-F-P139-042}{display}{none}
\[
  fx_*^{-3/2-\lambda},\qquad -fx_*^{-1/2-\lambda},\qquad \sqrt{2}x_*^{1/2}.
\]

\NSPage{140}%
\begin{EESegment}{EE-TGT-P140-001}
The three exponents are different, so Lemma~\ref{lem:A.1} says that this matrix is invertible. Its inverse is uniform
for \NSi{NS-F-P140-001}{\eta\in[-1,1]}. The first two rows have quadratic remainders, and the last row is linear.
The ratios \NSi{NS-F-P140-002}{X_*/X_K} and \NSi{NS-F-P140-003}{e_*/e_K} are fixed. After dividing
\eqref{eq:A.40}--\eqref{eq:A.42} by the scales in \eqref{eq:A.43}, the result is
\NSi{NS-F-P140-004}{O_m(X_K^{-1})}, even after any fixed number of
\NSi{NS-F-P140-005}{\eta}-derivatives. Lemma~\ref{lem:A.2} now gives a smooth local solution with
\NSi{NS-F-P140-006}{\lVert\partial_\eta^m c\rVert_\infty\leq C_mX_K^{-1}}. Choose
\NSi{NS-F-P140-007}{X_R} large enough that the moment equations have exactly one solution near zero correction
coefficients and that the Jacobian there is invertible. Implicit differentiation then gives a finite bound for every
fixed higher order of \NSi{NS-F-P140-008}{\eta}-derivatives. Only the finitely many derivative bounds used in the cone
argument have to meet specified smallness thresholds.
\end{EESegment}

\begin{EESegment}{EE-TGT-P140-002}
The three changed moments have now been restored. The functions \NSi{NS-F-P140-009}{M,J} never changed, because
\NSi{NS-F-P140-010}{U=0} on both edited regions. It remains to check that the local changes in the profiles and the
running integrals keep the stress inside the cone from the compensation patch to the terminal interval. The normalized
radial moment formulas \eqref{eq:4.16}, written in \NSi{NS-F-P140-011}{X/X_R}, carry the small changes in the profile
and its moments over to \NSi{NS-F-P140-012}{Q_s,N_s} on the fixed compact logarithmic interval from this patch to
\NSi{NS-F-P140-013}{y=.5}. More explicitly,
\NSi{NS-F-P140-014}{I=X_R^{3/2}\widehat I},
\NSi{NS-F-P140-015}{J=X_R^{3/2}\widehat J},
\NSi{NS-F-P140-016}{M=X_R\widehat M}, and
\NSi{NS-F-P140-017}{S=X_R\widehat S}. Substitute these into the formulas for
\NSi{NS-F-P140-018}{Q_s,N_s}. Every \NSi{NS-F-P140-019}{X_R}-factor then cancels, including the factors inside
\NSi{NS-F-P140-020}{H}. This comparison costs one more \NSi{NS-F-P140-021}{\eta}-derivative and gives
\NSi{NS-F-P140-022}{O(X_K^{-1})} changes in the needed \NSi{NS-F-P140-023}{\eta}-derivatives of
\NSi{NS-F-P140-024}{Q_s,N_s}. The positive lower bounds for \NSi{NS-F-P140-025}{E,Q_s}, the slope bounds, and the
admissible stress cone margins from Proposition~\ref{prop:A.4} persist for large \NSi{NS-F-P140-026}{X_R}. The
intermediate values of \NSi{NS-F-P140-027}{I} and the transition value of
\NSi{NS-F-P140-028}{Q_s} may change while those cone inequalities continue to hold.
\end{EESegment}

\begin{EESegment}{EE-TGT-P140-003}
Finally, the total difference in pressure is zero. Integrating
\NSi{NS-F-P140-029}{\Pi_X=E^2/(2X)} forward from the unchanged \NSi{NS-F-P140-030}{\Pi_0} still gives the
pressure normalized to vanish at radial infinity. Before the compensation patch, this forward solution is unchanged.
In particular, the terminal heat factor needs to be only smooth in \NSi{NS-F-P140-031}{\eta}; it does not change the
analytic pressure datum on the axis.
\end{EESegment}
\end{proof}

\begin{EESegment}{EE-TGT-P140-004}
\subsection{Finding the stress from the exterior profile.}\label{sec:A.7}
\end{EESegment}
\begin{EESegment}{EE-TGT-P140-005}
This subsection uses the corrected moments for two things. Lemma~\ref{lem:A.8} shows that the stress is zero outside
the annulus, and Proposition~\ref{prop:A.10} determines the stress direction as the outer edge is approached. First, the
moment identities will show that the full weighted integrals of the leading residual are zero. The stress can then be
found by integrating from the current radius out to infinity. Once Proposition~\ref{prop:B.2} and
Corollary~\ref{cor:B.10} have supplied the regular axis and the matching moment conditions, this formula depends only
on the exterior profile.
\end{EESegment}

\begin{lemma}\label{lem:A.8}
\begin{EESegment}{EE-TGT-P140-006}
Suppose the leading axisymmetric field is smooth at the axis. Suppose also that it has
\NSi{NS-F-P140-032}{U=0} and
\NSi{NS-F-P140-033}{E=E_{\mathrm{pow}}f_o[1+\chi_K(\mathcal H(2d/X)-1)]} for
\NSi{NS-F-P140-034}{y\geq 0}, and that it satisfies the exact moment conditions
\end{EESegment}
\NSFormula{NS-F-P140-035}{display}{none}
\[
  M(\infty)=J(\infty)=S(\infty)=0,\qquad
  \int_0^\infty(H-H_{\mathrm{pow}})\,dX=0,\qquad
  \Pi(X)=-\frac12\int_X^\infty\frac{E(x)^2}{x}\,dx.
\]
\begin{EESegment}{EE-TGT-P140-007}
Then the leading tangential stresses are zero for \NSi{NS-F-P140-036}{X\geq X_b}. On
\NSi{NS-F-P140-037}{y\geq 1/2}, where \NSi{NS-F-P140-038}{\chi_K=1}, equation~\eqref{eq:A.46} writes them as
weighted integrals of the leading residual from \NSi{NS-F-P140-039}{r} to infinity, with a positive sign in front of
each integral.
\end{EESegment}
\end{lemma}

\begin{proof}
\begin{EESegment}{EE-TGT-P140-008}
On the tail, the radial moment functions can be written as the following integrals from
\NSi{NS-F-P140-040}{X} to infinity. There \NSi{NS-F-P140-041}{U=M=J=0},
\NSi{NS-F-P140-042}{W=1}, and
\end{EESegment}
\NSFormula{NS-F-P140-043}{display}{none}
\[
  I(X)=\frac{XH_{\mathrm{pow}}(X)}{1-h}-\int_X^\infty(H-H_{\mathrm{pow}})\,dx,
  \qquad
  S(X)=\frac12\int_X^\infty E^2\,dx,
\]
\NSFormula{NS-F-P140-044}{display}{none}
\[
  Q_s=-1+\frac{(1-h)I-D\eta I_\eta}{XH},
  \qquad
  N_s=\frac{4h\eta S-dS_\eta}{X}+4A\eta\Pi-d\Pi_\eta.
\]

\NSPage{141}%
\begin{EESegment}{EE-TGT-P141-001}
The first formula uses \NSi{NS-F-P141-001}{h<1} to integrate
\NSi{NS-F-P141-002}{H_{\mathrm{pow}}} from zero. The remaining formulas follow directly from the exact moments and
the integral formulas \eqref{eq:4.16} for \NSi{NS-F-P141-003}{Q_s,N_s}.
\end{EESegment}

\begin{EESegment}{EE-TGT-P141-002}
The backward stress formula is valid only if no boundary term remains, so check the boundaries next. Let
\NSi{NS-F-P141-004}{K_{\mathrm{pow}}=c_\infty s^{-A}}. At large radius, uniformly for
\NSi{NS-F-P141-005}{\tau} in bounded nonnegative intervals,
\end{EESegment}
\NSFormula{NS-F-P141-006}{display}{A.44}
\begin{equation}
  K-K_{\mathrm{pow}}=O_h(\tau r^{-3-2h}),\qquad
  \partial_tK=O_h(r^{-3-2h}),\qquad
  r^2K_r-rK=O_h(r^{-2h}).
  \tag{A.44}\label{eq:A.44}
\end{equation}
\begin{EESegment}{EE-TGT-P141-003}
These bounds come from \eqref{eq:A.36} and the corresponding bounds for any fixed derivative. In particular, at any
sufficiently large fixed physical radius \NSi{NS-F-P141-007}{R_*},
\end{EESegment}
\NSFormula{NS-F-P141-008}{display}{none}
\[
  \int_{R_*}^\infty\bigl(r^2|K-K_{\mathrm{pow}}|+|\partial_tK|\bigr)\,dr
  \leq C_h(1+\tau)\int_{R_*}^\infty r^{-1-2h}\,dr
  =\frac{C_h(1+\tau)}{2h}R_*^{-2h}.
\]
\begin{EESegment}{EE-TGT-P141-004}
So the subtracted angular moment converges. The same integrable upper bound works uniformly on compact time intervals,
which justifies differentiating that moment with respect to time. This gives
\end{EESegment}
\NSFormula{NS-F-P141-009}{display}{A.45}
\begin{equation}
  \int_0^\infty r^2(u_\theta-K_{\mathrm{pow}})\,dr
  =q^{3/2-A}\int_0^\infty(H-H_{\mathrm{pow}})\,dX=0.
  \tag{A.45}\label{eq:A.45}
\end{equation}
\begin{EESegment}{EE-TGT-P141-005}
At the axis, the subtracted power has weight \NSi{NS-F-P141-010}{r^{1-2h}}, which is also integrable. Its time
derivative is zero. The angular transport integral is the sum of the radial boundary
\NSi{NS-F-P141-011}{[r^2u_ru_\theta]_0^\infty} and the axial derivative of
\NSi{NS-F-P141-012}{\int r^2u_zu_\theta\,dr=q^{3/2-2A}J(\infty)=0}. Axis regularity makes the contribution at the axis
zero, while \NSi{NS-F-P141-013}{U=M=0} on the exterior makes the contribution at infinity zero. The latter also gives
\NSi{NS-F-P141-014}{V_0^{(0)}=0}. By \eqref{eq:A.44}, radial viscosity integrates to
\NSi{NS-F-P141-015}{[r^2\partial_ru_\theta-ru_\theta]_0^\infty=0}. The leading angular residual
\NSi{NS-F-P141-016}{\mathscr R_\theta^{(0)}} therefore has zero integral against
\NSi{NS-F-P141-017}{r^2\,dr}.
\end{EESegment}

\begin{EESegment}{EE-TGT-P141-006}
For the axial component, \NSi{NS-F-P141-018}{\int ru_z\,dr=q^{1-A}M(\infty)=0}. The pressure formula
\NSi{NS-F-P141-019}{p(r)=-\int_r^\infty u_\theta(r')^2\,dr'/r'} gives
\end{EESegment}
\NSFormula{NS-F-P141-020}{display}{none}
\[
  \int_0^\infty r(u_z^2+p)\,dr
  =q^{1-2A}\int_0^\infty(U^2-E^2/2)\,dX=0.
\]
\begin{EESegment}{EE-TGT-P141-007}
The exact pressure calculation behind this identity is
\end{EESegment}
\NSFormula{NS-F-P141-021}{display}{none}
\[
  \int_0^\infty rp(r)\,dr
  =-\int_0^\infty\frac{u_\theta(r')^2}{r'}
    \left(\int_0^{r'}r\,dr\right)dr'
  =-\frac12\int_0^\infty r'u_\theta(r')^2\,dr'.
\]
\begin{EESegment}{EE-TGT-P141-008}
The exterior bound \NSi{NS-F-P141-022}{p=O(r^{-2-4h})} makes the integral converge and gives
\NSi{NS-F-P141-023}{r^2p\to 0}. The boundary terms from radial transport and viscosity in the axial equation are also
zero. So \NSi{NS-F-P141-024}{\int r\mathscr R_z^{(0)}\,dr=0}.
\end{EESegment}

\begin{EESegment}{EE-TGT-P141-009}
Let \NSi{NS-F-P141-025}{T_\theta,T_z} be the physical stress components, so
\NSi{NS-F-P141-026}{T=q^{-A-1/2}\mathcal T_0}. Their defining integrals from the axis can now be written as
\end{EESegment}
\NSFormula{NS-F-P141-027}{display}{A.46}
\begin{equation}
  T_\theta(r)=\frac1{r^2}\int_r^\infty r'^2\mathscr R_\theta^{(0)}(r')\,dr',
  \qquad
  T_z(r)=\frac1r\int_r^\infty r'\mathscr R_z^{(0)}(r')\,dr'.
  \tag{A.46}\label{eq:A.46}
\end{equation}
\begin{EESegment}{EE-TGT-P141-010}
For \NSi{NS-F-P141-028}{X\geq X_b}, the field is
\NSi{NS-F-P141-029}{(u_r,u_\theta,u_z)=(0,K,0)}, and its pressure does not depend on
\NSi{NS-F-P141-030}{z} and is normalized to vanish at radial infinity. Lemma~\ref{lem:A.6} says that both leading
residuals are zero there, so the two stress components are zero for \NSi{NS-F-P141-031}{X\geq X_b}.
\end{EESegment}
\end{proof}

\begin{EESegment}{EE-TGT-P141-011}
The backward integrals contain the cutoff factor, which goes to zero to every order at the outer edge. The next lemma
pulls this factor away from a smooth coefficient. This makes it possible to estimate the direction of the stress even
as the stress itself goes to zero.
\end{EESegment}

\begin{lemma}\label{lem:A.9}
\begin{EESegment}{EE-TGT-P141-012}
Let \NSi{NS-F-P141-032}{c>0}, let \NSi{NS-F-P141-033}{j} be a fixed nonnegative integer, and let
\NSi{NS-F-P141-034}{b(u,\eta)} be smooth for \NSi{NS-F-P141-035}{0\leq u\leq\delta_0} and for
\NSi{NS-F-P141-036}{\eta} in a compact interval \NSi{NS-F-P141-037}{K}. Then there is a coefficient
\NSi{NS-F-P141-038}{B} such that, whenever \NSi{NS-F-P141-039}{0<\delta\leq\delta_0},
\end{EESegment}
\NSFormula{NS-F-P141-040}{display}{A.47}
\begin{equation}
  \int_0^\delta e^{-c/u^2}u^{-j}b(u,\eta)\,du
  =\frac12e^{-c/\delta^2}\delta^{3-j}B(\delta,\eta).
  \tag{A.47}\label{eq:A.47}
\end{equation}

\NSPage{142}%
\begin{EESegment}{EE-TGT-P142-001}
The coefficient \NSi{NS-F-P142-001}{B} is smooth on
\NSi{NS-F-P142-002}{[0,\delta_0]\times K}. At the endpoint,
\NSi{NS-F-P142-003}{B(0,\eta)=b(0,\eta)/c}, and every fixed mixed derivative is bounded there. The same result holds
when \NSi{NS-F-P142-004}{\eta} is replaced by any finite collection of smooth compact parameters.
\end{EESegment}
\end{lemma}

\begin{proof}
\begin{EESegment}{EE-TGT-P142-002}
Make the substitution \NSi{NS-F-P142-005}{u=\delta/(1+\delta^2v)^{1/2}}. Its Jacobian is
\NSi{NS-F-P142-006}{-\tfrac12\delta^3(1+\delta^2v)^{-3/2}}, and
\NSi{NS-F-P142-007}{c/u^2=c/\delta^2+cv}. This proves \eqref{eq:A.47} with
\end{EESegment}
\NSFormula{NS-F-P142-008}{display}{none}
\[
  B(\delta,\eta)=\int_0^\infty e^{-cv}(1+\delta^2v)^{(j-3)/2}
  b\!\left(\frac{\delta}{\sqrt{1+\delta^2v}},\eta\right)\,dv.
\]
\begin{EESegment}{EE-TGT-P142-003}
Every fixed mixed derivative of this last integrand is bounded by \NSi{NS-F-P142-009}{e^{-cv}} times a polynomial in
\NSi{NS-F-P142-010}{v}, with one bound for the whole closed rectangle of parameters. Differentiation under the
integral is valid and proves both smoothness and the derivative bounds. Setting
\NSi{NS-F-P142-011}{\delta=0} inside the integral gives the stated value at the endpoint.
\end{EESegment}
\end{proof}

\begin{EESegment}{EE-TGT-P142-004}
The stress is zero at the outer edge, so its limiting direction depends on how fast its two components go to zero
relative to each other. The backward formula and the flat-integral lemma give both rates. They also give a positive
leading coefficient for the angular component.
\end{EESegment}

\begin{proposition}\label{prop:A.10}
\begin{EESegment}{EE-TGT-P142-005}
Under the regular-axis and exact-moment assumptions of Lemma~\ref{lem:A.8}, the profile on the terminal interval stays
inside the admissible stress cone for \NSi{NS-F-P142-012}{.5\leq y<3}. The direction of the stress extends smoothly
to the outer endpoint. Through this extension,
\NSi{NS-F-P142-013}{2-(a-2)(T_z/T_\theta)^2} has a uniform positive lower bound. More exactly, when
\NSi{NS-F-P142-014}{\delta=3-y} is small enough,
\end{EESegment}
\NSFormula{NS-F-P142-015}{display}{A.48}
\begin{equation}
  \mathcal T_{0,\theta}=e^{-4/\delta^2}\delta^{-3}b_\theta(\delta,\eta),
  \qquad b_\theta(0,\eta)>0,
  \tag{A.48}\label{eq:A.48}
\end{equation}
\NSFormula{NS-F-P142-016}{display}{A.49}
\begin{equation}
  \mathcal T_{0,z}=e^{-4/\delta^2}\delta^3b_z(\delta,\eta),
  \tag{A.49}\label{eq:A.49}
\end{equation}
\NSFormula{NS-F-P142-017}{display}{A.50}
\begin{equation}
  \frac{\mathcal T_{0,z}}{\mathcal T_{0,\theta}}=\delta^6\frac{b_z}{b_\theta}\longrightarrow 0.
  \tag{A.50}\label{eq:A.50}
\end{equation}
\begin{EESegment}{EE-TGT-P142-006}
Both coefficients are smooth on a closed collar around the outer endpoint, and
\NSi{NS-F-P142-018}{b_\theta} has a positive lower bound there. For every fixed mixed profile derivative
\NSi{NS-F-P142-019}{\partial^I}, there are finite constants \NSi{NS-F-P142-020}{C_I,N_I} such that
\end{EESegment}
\NSFormula{NS-F-P142-021}{display}{A.51}
\begin{equation}
  |\partial^I\mathcal T_0|\leq C_Ie^{-4/\delta^2}\delta^{-N_I},
  \qquad |\mathcal T_0|\geq ce^{-4/\delta^2}\delta^{-3}.
  \tag{A.51}\label{eq:A.51}
\end{equation}
\begin{EESegment}{EE-TGT-P142-007}
These constants may depend on all the fixed choices in the profile and on \NSi{NS-F-P142-022}{I}, but they do not
depend on the physical scale \NSi{NS-F-P142-023}{q}.
\end{EESegment}
\end{proposition}

\begin{proof}
\begin{EESegment}{EE-TGT-P142-008}
First find a positive angular stress and bound the axial stress relative to it. Then factor out the terms that become
flat at the endpoint, which proves that the limiting direction is smooth.
\end{EESegment}

\begin{EESegment}{EE-TGT-P142-009}
On \NSi{NS-F-P142-024}{y\geq 1/2}, where \NSi{NS-F-P142-025}{\chi_K=1}, let
\NSi{NS-F-P142-026}{f=f_o(y)} and write \NSi{NS-F-P142-027}{f'=\partial_yf}. Here
\NSi{NS-F-P142-028}{u_\theta=Kf} and \NSi{NS-F-P142-029}{u_r=u_z=0}. Use the heat equation from
Lemma~\ref{lem:A.6}, together with \NSi{NS-F-P142-030}{y_t=(qL)^{-1}} and
\NSi{NS-F-P142-031}{y_z=-2\eta/(q^DL)}. This gives
\end{EESegment}
\NSFormula{NS-F-P142-032}{display}{A.52}
\begin{equation}
  \mathscr R_\theta^{(0)}
  =\frac{Kf'}{qL}-K(f_{rr}+r^{-1}f_r)-2K_rf_r,
  \tag{A.52}\label{eq:A.52}
\end{equation}
\NSFormula{NS-F-P142-033}{display}{A.53}
\begin{equation}
  p_z(r)=\frac{2\eta}{q^DL}\int_y^3K(y')^2f(y')f'(y')\,dy'.
  \tag{A.53}\label{eq:A.53}
\end{equation}
\begin{EESegment}{EE-TGT-P142-010}
Integrate the viscous terms in \eqref{eq:A.52} by parts inside \eqref{eq:A.46}. The exact result is
\end{EESegment}
\NSFormula{NS-F-P142-034}{display}{A.54}
\begin{equation}
  T_\theta
  =Kf_r+\frac1{r^2}\int_r^\infty(r'K-r'^2K_{r'})f_{r'}\,dr'
  +\frac1{qLr^2}\int_r^\infty r'^2Kf'\,dr'.
  \tag{A.54}\label{eq:A.54}
\end{equation}
\begin{EESegment}{EE-TGT-P142-011}
All three terms are nonnegative, because \NSi{NS-F-P142-035}{K>0},
\NSi{NS-F-P142-036}{K_r<0}, and \NSi{NS-F-P142-037}{f'\geq 0}. Along the rest of the logarithmic interval, whose
length is at most \NSi{NS-F-P142-038}{2.5}, both the ratios of radii and the ratios of heat factors have upper and
lower bounds.
\end{EESegment}

\NSPage{143}%
\begin{EESegment}{EE-TGT-P143-001}
Since \NSi{NS-F-P143-001}{\int_y^3f'(v)\,dv=\rho_o\psi_o(y)}, the last term in \eqref{eq:A.54} gives
\end{EESegment}
\NSFormula{NS-F-P143-002}{display}{none}
\[
  T_\theta\geq c\frac{rK}{qL}\rho_o\psi_o(y)>0\qquad(y<3).
\]
\begin{EESegment}{EE-TGT-P143-002}
In the same way, \eqref{eq:A.53} and the integral formula for \NSi{NS-F-P143-003}{T_z} give
\end{EESegment}
\NSFormula{NS-F-P143-004}{display}{A.55}
\begin{equation}
  |p_z|\leq Cq^{-D}K^2\rho_o\psi_o,
  \qquad |T_z|\leq Crq^{-D}K^2\rho_o\psi_o,
  \qquad \left|\frac{T_z}{T_\theta}\right|\leq Cq^{1-D}K
  =CE_{\mathrm{pow}}\mathcal H(2d/X).
  \tag{A.55}\label{eq:A.55}
\end{equation}
\begin{EESegment}{EE-TGT-P143-003}
The constants here are uniform when \NSi{NS-F-P143-005}{X_R} is large enough and
\NSi{NS-F-P143-006}{h} is small. The last equality uses \NSi{NS-F-P143-007}{1-D=A}. The drop in amplitude during
the exterior transition in \eqref{eq:A.17} makes the final expression small. At the same time,
\NSi{NS-F-P143-008}{b_s=0}, and
\end{EESegment}
\NSFormula{NS-F-P143-009}{display}{A.56}
\begin{equation}
  a=2+2h+2Z\frac{\mathcal H'}{\mathcal H}-2\frac{f'}f>2+h,
  \qquad Z=2d/X,
  \tag{A.56}\label{eq:A.56}
\end{equation}
\begin{EESegment}{EE-TGT-P143-004}
because \NSi{NS-F-P143-010}{f'/f<h/4} and, once \NSi{NS-F-P143-011}{X_R} is large enough,
\NSi{NS-F-P143-012}{-Z\mathcal H'/\mathcal H<h/4}. Also, \NSi{NS-F-P143-013}{a\leq 2+2h}. It follows that
\NSi{NS-F-P143-014}{v_s=a},
\NSi{NS-F-P143-015}{P_{\mathrm c}-v_s=\mathcal T_{0,\theta}/F>0}, and
\NSi{NS-F-P143-016}{J_{\mathrm c}=\mathcal T_{0,z}/F}. The directional cone inequality is
\NSi{NS-F-P143-017}{(a-2)(T_z/T_\theta)^2<2}. Equation~\eqref{eq:A.55} leaves a fixed positive margin in this
inequality. This proves the admissible stress cone at every point before the outer endpoint
\NSi{NS-F-P143-018}{X=X_b}.
\end{EESegment}

\begin{EESegment}{EE-TGT-P143-005}
The cone inequality now holds at every interior point of the terminal interval. It remains to show that the direction
of the stress extends smoothly to the endpoint. Factor out the part shared by the two stress components as they go to
zero, and then show that the angular coefficient stays positive. Near \NSi{NS-F-P143-019}{\delta=0}, the chosen
smooth step has the form
\end{EESegment}
\NSFormula{NS-F-P143-020}{display}{none}
\[
  \psi_o=e^{-4/\delta^2}g(\delta),
  \qquad
  g(\delta)=\frac{1}{e^{-(1-\delta/2)^{-2}}+e^{-4/\delta^2}},
  \qquad g(0)=e,
\]
\begin{EESegment}{EE-TGT-P143-006}
and so
\end{EESegment}
\NSFormula{NS-F-P143-021}{display}{none}
\[
  f'=\rho_oe^{-4/\delta^2}\delta^{-3}\bigl[8g(\delta)+\delta^3g'(\delta)\bigr].
\]
\begin{EESegment}{EE-TGT-P143-007}
Apply Lemma~\ref{lem:A.9} to the backward formulas for the stress.
\end{EESegment}

\begin{EESegment}{EE-TGT-P143-008}
Multiply the boundary term in \eqref{eq:A.54} by \NSi{NS-F-P143-022}{q^{A+1/2}}. This gives
\end{EESegment}
\NSFormula{NS-F-P143-023}{display}{none}
\[
  q^{A+1/2}Kf_r=\frac{2E_{\mathrm{pow}}(X)\mathcal H(2d/X)}{\sqrt{2X}}f'.
\]
\begin{EESegment}{EE-TGT-P143-009}
Both integral terms contain \NSi{NS-F-P143-024}{f'}. Equation~\eqref{eq:A.47}, with
\NSi{NS-F-P143-025}{j=3,c=4}, pulls out a factor \NSi{NS-F-P143-026}{e^{-4/\delta^2}} and leaves a smooth
coefficient. Compared with the boundary factor, those two terms vanish to at least order
\NSi{NS-F-P143-027}{\delta^3}. This proves \eqref{eq:A.48}, with
\end{EESegment}
\NSFormula{NS-F-P143-028}{display}{none}
\[
  b_\theta(0,\eta)=\frac{16\rho_oE_{\mathrm{pow}}(X_b)}{\sqrt{2X_b}}
  \mathcal H(2d/X_b)g(0)>0.
\]
\begin{EESegment}{EE-TGT-P143-010}
The positive lower bound is uniform for \NSi{NS-F-P143-029}{-1\leq\eta\leq 1}, because
\NSi{NS-F-P143-030}{\mathcal H} has a positive minimum on \NSi{NS-F-P143-031}{[0,2/X_b]}.
\end{EESegment}

\begin{EESegment}{EE-TGT-P143-011}
The pressure integral \eqref{eq:A.53} contains one \NSi{NS-F-P143-032}{f'} factor. The same identity first writes
\NSi{NS-F-P143-033}{p_z=e^{-4/\delta^2}} times a smooth coefficient, apart from the factor
\NSi{NS-F-P143-034}{q^{-D-2A}}. The integral formula for \NSi{NS-F-P143-035}{T_z} requires one more integration with
smooth weights. Applying \eqref{eq:A.47} with \NSi{NS-F-P143-036}{j=0} gives exactly \eqref{eq:A.49}. In this
normalization, the powers cancel because
\NSi{NS-F-P143-037}{q^{A+1/2}q^{1/2-D-2A}=q^{1-D-A}=1}. Every coefficient left over contains only smooth
functions of \NSi{NS-F-P143-038}{X,\eta,\mathcal H(2d/X)} and \NSi{NS-F-P143-039}{L^{-1}}; here
\NSi{NS-F-P143-040}{L\geq 1-2h>0}. This also proves smoothness on the closed parameter ranges and shows that neither
flat factor contains the physical scale.
\end{EESegment}

\NSPage{144}%
\begin{EESegment}{EE-TGT-P144-001}
Divide by the positive coefficient \NSi{NS-F-P144-001}{b_\theta} to get \eqref{eq:A.50}. The limiting direction is
\NSi{NS-F-P144-002}{(1,0)}, where the normalized directional quadratic expression equals
\NSi{NS-F-P144-003}{2}. The stress itself is zero at the endpoint. Finally, every fixed derivative of the profile
factorizations adds only a finite inverse power of \NSi{NS-F-P144-004}{\delta}. This proves \eqref{eq:A.51}. A smooth
positive factor from the other endpoint may multiply this outer weight without changing any of these local conclusions.
\end{EESegment}
\end{proof}

\begin{EESegment}{EE-TGT-P144-002}
\section{Analytic profiles near the axis and how to continue them}\label{sec:B}
\end{EESegment}

\begin{EESegment}{EE-TGT-P144-003}
This appendix builds the inner part of the leading profile \NSi{NS-F-P144-005}{(U,E,\Pi)}. On this inner part, the
leading residual stress \NSi{NS-F-P144-006}{\mathcal T_0} must be zero, as Proposition~\ref{prop:B.2} requires. The
outer profile has already fixed the pressure datum \NSi{NS-F-P144-007}{\Pi_0} on the axis, in
Lemma~\ref{lem:A.5}. With that datum fixed, choose the remaining axis data
\NSi{NS-F-P144-008}{U_*,\phi_*} and solve the nonlinear profile equations \eqref{eq:4.13} on a short radial interval.
The solution must be regular at the axis. At its outer endpoint, it must also satisfy the strict shear inequality
\eqref{eq:B.19}, so that Corollary~\ref{cor:B.10} can continue it toward the outer profile. The large parameter
\NSi{NS-F-P144-009}{\Lambda} in the azimuthal datum \eqref{eq:B.3} gives both the analytic solution and the endpoint
bounds needed for this continuation, as Propositions~\ref{prop:B.2} and~\ref{prop:B.3} show. Equations~\eqref{eq:4.4}--
\eqref{eq:4.5} then show that the corresponding velocity is smooth in Cartesian coordinates across the axis.
\end{EESegment}

\begin{EESegment}{EE-TGT-P144-004}
\subsection{Choosing the axis data.}\label{sec:B.1}
\end{EESegment}
\begin{EESegment}{EE-TGT-P144-005}
The pressure datum \NSi{NS-F-P144-010}{\Pi_0} is already prescribed. The azimuthal datum
\NSi{NS-F-P144-011}{\phi_*=\phi(0,\eta)} and the axial velocity \NSi{NS-F-P144-012}{U_*} on the axis still have to
be chosen. Choose them here so that the angular source \NSi{NS-F-P144-013}{S_q} stays positive and the endpoint
inequality holds throughout the parameter interval, as Proposition~\ref{prop:B.3} requires. Equation~\eqref{eq:B.2}
separates the interval into two parts; the azimuthal term gives the needed lower bound on one part, and the axial term
gives it on the other.
\end{EESegment}

\begin{EESegment}{EE-TGT-P144-006}
Use the datum \NSi{NS-F-P144-014}{\Pi_0} from Lemma~\ref{lem:A.5}. It is real analytic near
\NSi{NS-F-P144-015}{[-1,1]}, it is even, and it satisfies
\end{EESegment}
\NSFormula{NS-F-P144-016}{display}{none}
\[
  \Pi_0\leq-cP_*^2f^2,
  \qquad \eta\Pi_{0\eta}>0\quad(\eta\neq 0),
  \qquad f=(1+\eta^2)^{-1}.
\]
\begin{EESegment}{EE-TGT-P144-007}
Every schedule parameter, including \NSi{NS-F-P144-017}{0<h\leq 10^{-2}}, is already fixed. Choose
\NSi{NS-F-P144-018}{0<j_0\leq .05} small, and set
\end{EESegment}
\NSFormula{NS-F-P144-019}{display}{B.1}
\begin{equation}
\begin{aligned}
  U_*&=4\eta+j_0, & H_*&=D\eta+dU_*,\\
  W_*&=1-dU_{*\eta}-2D\eta U_*,
  & Z_*&=-A(1-2\eta U_*)U_*-H_*U_{*\eta}-d\Pi_{0\eta}+4A\eta\Pi_0.
\end{aligned}
\tag{B.1}\label{eq:B.1}
\end{equation}
\begin{EESegment}{EE-TGT-P144-008}
The function \NSi{NS-F-P144-020}{H_*} has exactly one zero
\NSi{NS-F-P144-021}{\eta_0\in(-1,0)}. To see why, note that on \NSi{NS-F-P144-022}{(-1,1)} the function
\NSi{NS-F-P144-023}{H_*/d=D\eta/d+4\eta+j_0} is strictly increasing from
\NSi{NS-F-P144-024}{-\infty} to \NSi{NS-F-P144-025}{+\infty}. Its zero satisfies
\NSi{NS-F-P144-026}{|\eta_0|\asymp j_0}. Also,
\end{EESegment}
\NSFormula{NS-F-P144-027}{display}{none}
\[
  -W_*=3-8h\eta^2+(1-2h)j_0\eta>2.8,
  \qquad Z_*(\eta_0)\geq cj_0P_*^2>0.
\]
\begin{EESegment}{EE-TGT-P144-009}
For the second inequality, \NSi{NS-F-P144-028}{4A\eta_0\Pi_0(\eta_0)} is positive and at least
\NSi{NS-F-P144-029}{cj_0P_*^2}. The term \NSi{NS-F-P144-030}{-d\Pi_{0\eta}(\eta_0)} is nonnegative, the
\NSi{NS-F-P144-031}{H_*} term is zero, and the term that remains is \NSi{NS-F-P144-032}{O(j_0)}. The choice of
\NSi{NS-F-P144-033}{P_*}, which is already large, absorbs this last term.
\end{EESegment}

\begin{EESegment}{EE-TGT-P144-010}
It is now possible to choose \NSi{NS-F-P144-034}{\delta_*>0} first, so that the set
\NSi{NS-F-P144-035}{\{|Z_*|\leq\delta_*\}} stays away from a neighborhood of
\NSi{NS-F-P144-036}{\eta_0}. Then choose \NSi{NS-F-P144-037}{\sigma_*>0} small enough that
\end{EESegment}
\NSFormula{NS-F-P144-038}{display}{B.2}
\begin{equation}
  \chi=\frac{H_*^2}{H_*^2+\sigma_*^2}>.99
  \qquad\text{where }|Z_*|\leq\delta_*.
  \tag{B.2}\label{eq:B.2}
\end{equation}
\begin{EESegment}{EE-TGT-P144-011}
Make both choices on a slightly larger compact interval \NSi{NS-F-P144-039}{I\supset[-1,1]}. Then choose a bounded,
simply connected complex neighborhood \NSi{NS-F-P144-040}{\Omega} of \NSi{NS-F-P144-041}{I}. On the closure of
this neighborhood, neither \NSi{NS-F-P144-042}{L} nor \NSi{NS-F-P144-043}{H_*^2+\sigma_*^2} may vanish.
\end{EESegment}
